\documentclass[11pt]{article}
\usepackage[T1]{fontenc}
\usepackage{textcomp}
\usepackage[margin=0.75in]{geometry}
\usepackage{amsmath,amssymb,amsthm}
\usepackage{array,booktabs}
\usepackage{hyperref}
\setlength{\parindent}{0pt}
\newtheorem{theorem}{Theorem}
\theoremstyle{definition}
\newtheorem{definition}{Definition}
\title{Flow Proof Artifact: prop-15-part-to-part-as-whole}
\author{Generated by \texttt{flow doc proof}}
\date{Flow Proof Book}
\begin{document}
\maketitle
A part has to a part the same ratio which the corresponding whole has to the corresponding whole.\\[0.5em]
\textbf{Source.} euclid — Elements, Book V, Proposition 15\\[0.5em]
\bigskip
\noindent{\large \textbf{Derived fact 1.} \textit{A part has to a part the same ratio which the corresponding whole has to the corresponding whole} \hfill the Euclidean plane \cdot Euclid Book V \cdot Proposition 15: a part has the same ratio to its part as the whole to its whole}\par
\smallskip\noindent\textit{Coordinate: the Euclidean plane \cdot Euclid Book V \cdot Proposition 15: a part has the same ratio to its part as the whole to its whole}\par
\smallskip\noindent\textit{Source: euclid — Elements, Book V, Proposition 15}\par
\smallskip\noindent\textit{Built on: proposition 6: magnitudes with the same ratio to a third are proportional to one another, for Euclid Book V on the Euclidean plane}\par
\medskip
\noindent\textbf{Goal.} A part has to a part the same ratio which the corresponding whole has to the corresponding whole\par
\noindent$\displaystyle \text{part ratio equals whole ratio}$\par
\medskip
\noindent
\begin{tabular}{@{} >{\raggedright\arraybackslash}p{0.06\textwidth} >{\raggedright\arraybackslash}p{0.47\textwidth} >{\raggedleft\arraybackslash}p{0.41\textwidth} @{}}
\textbf{\#} & \textbf{Proof} & \textbf{Mathematics} \\
\midrule
\textbf{1.} & We prove that proposition 15: a part has the same ratio to its part as the whole to its whole for Euclid Book V the Euclidean plane. & \\[0.45em]
\textbf{2.} & Let part\_a = part\_of\_first. & \\[0.45em]
\textbf{3.} & Let part\_b = part\_of\_second. & \\[0.45em]
\textbf{4.} & Let whole\_a = whole\_first. & \\[0.45em]
\textbf{5.} & From step 2, step 3, and step 4, we can deduce that ratio(part a, part b) equals ratio(whole a, whole b). & $\displaystyle ratio(\text{part a}, \text{part b}) = ratio(\text{whole a}, \text{whole b})$ \\[0.45em]
\textbf{6.} & We invoke the derived fact governing Euclid Book V the Euclidean plane: proposition 6: magnitudes with the same ratio to a third are proportional to one another, for Euclid Book V on the Euclidean plane. & \\[0.45em]
\textbf{7.} & From step 6, this implies part ratio equals whole ratio. Hence proven. & $\displaystyle \text{part ratio equals whole ratio}$ \\[0.45em]
\bottomrule
\end{tabular}
\label{thm:Geometry-Euclid Book V-proposition_15_a_part_has_the_same_ratio_to_its_part_as_the_whole_to_its_whole}
\par\medskip\hrule
\end{document}
